% !TEX encoding = UTF-8
% !TEX program = latexmk
\documentclass{CWNUExam}

\cwnuexamsetup{
pdf = {
title    = {计算机学院软件工程专业2024级JAVA编程语言期末考试试卷A卷},
author   = {Computer Science Department},
subject  = {For internal academic use only},
keywords = {Java; Programming; Final Exam},
},
}

\cwnuexamsetup{
title = {
year       = {\the\year},
semester   = {1},
course     = {JAVA Programming Language},
suffix     = {期末},
type       = {A},
total-part = {4},
college    = {计算机学院},
major      = {软件工程专业},
grade      = {2024级},
exam-type  = {闭卷},
exam-time  = {120},
exam-category = {考试},
test-type  = {normal},
},
}

\examsetup{page/size = a3paper}

\begin{document}
\maketitle

\section{Section A: Multiple Choice Questions (15 \times 3 = 45)}

\begin{question}
Which keyword is used to inherit a class in Java?
\paren
\begin{choices}
\item implements
\item extends
\item new
\item super
\end{choices}
\end{question}

\begin{question}
Which data type is used to store a single character?
\paren
\begin{choices}
\item String
\item char
\item int
\item boolean
\end{choices}
\end{question}

\begin{question}
Which of the following is NOT a primitive data type?
\paren
\begin{choices}
\item int
\item double
\item boolean
\item String
\end{choices}
\end{question}

\begin{question}
Which method is the entry point of a Java program?
\paren
\begin{choices}
\item start()
\item run()
\item main()
\item init()
\end{choices}
\end{question}

\begin{question}
Which access modifier allows access only inside the same class?
\paren
\begin{choices}
\item public
\item protected
\item private
\item default
\end{choices}
\end{question}

\begin{question}
Which keyword is used to create an object in Java?
\paren
\begin{choices}
\item class
\item object
\item new
\item this
\end{choices}
\end{question}

\begin{question}
Which keyword refers to the current object?
\paren
\begin{choices}
\item super
\item this
\item current
\item self
\end{choices}
\end{question}

\begin{question}
Which loop runs at least one time?
\paren
\begin{choices}
\item for
\item while
\item do-while
\item foreach
\end{choices}
\end{question}

\begin{question}
Which keyword is used to manually throw an exception?
\paren
\begin{choices}
\item try
\item catch
\item throw
\item throws
\end{choices}
\end{question}

\begin{question}
Which block always executes whether exception occurs or not?
\paren
\begin{choices}
\item try
\item catch
\item final
\item finally
\end{choices}
\end{question}

\begin{question}
Which feature supports method overriding?
\paren
\begin{choices}
\item Encapsulation
\item Inheritance
\item Abstraction
\item Compilation
\end{choices}
\end{question}

\begin{question}
Which of the following best describes polymorphism?
\paren
\begin{choices}
\item Same method name, different classes
\item Same variable name
\item Same class name
\item Same package name
\end{choices}
\end{question}

\begin{question}
Which of the following is an interface keyword?
\paren
\begin{choices}
\item class
\item interface
\item abstract
\item implements
\end{choices}
\end{question}

\begin{question}
Which exception occurs when dividing by zero?
\paren
\begin{choices}
\item IOException
\item NullPointerException
\item ArithmeticException
\item ArrayException
\end{choices}
\end{question}

\begin{question}
Which collection is used to store elements in order?
\paren
\begin{choices}
\item Set
\item Map
\item List
\item Queue
\end{choices}
\end{question}


\section{Section B: Fill in the Blanks (5 \times 1 = 5)}

\begin{question}
Java programs start execution from the \fillin method.
\end{question}

\begin{question}
The parent class of all Java classes is \fillin.
\end{question}

\begin{question}
The keyword used to inherit a class is \fillin.
\end{question}

\begin{question}
The \fillin block is used to handle exceptions.
\end{question}

\begin{question}
Data hiding is achieved using \fillin.
\end{question}


% =========================
% Part C
% =========================
\section{Section C: Explain / Short Answer Questions (5 \times 2 = 10)}

\begin{problem}
What is Encapsulation in Java? (Answer in 2--3 sentences)
\end{problem}

\begin{problem}
What is Inheritance? Mention one advantage.
\end{problem}

\begin{problem}
What is method overloading?
\end{problem}

\begin{problem}
What is an interface in Java?
\end{problem}

\begin{problem}
What is exception handling?
\end{problem}

% =========================
% Part D
% =========================
\section{Section D: Programming / Design Questions (2 \times 20 = 40)}

\begin{problem}
    \textbf{Design Question 1 (Student System using Inheritance)}
    
    Design a simple Student system using inheritance.
    
    \textbf{Guidelines:}
    \begin{enumerate}
      \item Create a base class \texttt{Student} with: \texttt{name}, \texttt{rollNo}.
      \item Create a subclass \texttt{Result} that stores marks of 3 subjects.
      \item Add a method to calculate total marks.
      \item Display student details and total marks.
    \end{enumerate}
    
    \textbf{Note:}
    \begin{itemize}
      \item Full Java syntax is NOT required.
      \item Class structure + key methods are enough.
    \end{itemize}
    \end{problem}
    

    \begin{problem}
        \textbf{Design Question 2 (Bank System using Interface)}
        
        Design a simple Bank system.
        
        \textbf{Guidelines:}
        \begin{enumerate}
          \item Create an interface \texttt{Bank} with methods \texttt{deposit()} and \texttt{withdraw()}.
          \item Create a class \texttt{MyBank} that implements the interface.
          \item Use \texttt{try-catch} to handle withdrawal when balance is insufficient.
        \end{enumerate}
        
        \textbf{Note:}
        \begin{itemize}
          \item Pseudo-code or partial code is acceptable.
          \item Focus on logic, not full syntax.
        \end{itemize}
        \end{problem}
        

% =========================
% Answer Sheet (One Page)
% =========================
\newpage

\noindent\textbf{Student Name:} \underline{\hspace{6cm}}

\textbf{Student ID:} \underline{\hspace{6cm}}

\vspace{0.6cm}
\noindent\textbf{Section A: Multiple Choice (1--15)}\\[0.2cm]
\noindent
\begin{tabular}{|c|c|c|c|c|}
\hline
1 & 2 & 3 & 4 & 5  \\
\hline
\underline{\hspace{1.6cm}} & \underline{\hspace{1.6cm}} & \underline{\hspace{1.6cm}} & \underline{\hspace{1.6cm}} & \underline{\hspace{1.6cm}}  \\
\hline
\end{tabular}

\begin{tabular}{|c|c|c|c|c|}
\hline
6 & 7 & 8 & 9 & 10  \\
\hline
\underline{\hspace{1.6cm}} & \underline{\hspace{1.6cm}} & \underline{\hspace{1.6cm}} & \underline{\hspace{1.6cm}} & \underline{\hspace{1.6cm}}  \\
\hline
\end{tabular}

\begin{tabular}{|c|c|c|c|c|}
\hline
11 & 12 & 13 & 14 & 15 \\
\hline
\underline{\hspace{1.6cm}} & \underline{\hspace{1.6cm}} & \underline{\hspace{1.6cm}} & \underline{\hspace{1.6cm}} & \underline{\hspace{1.6cm}}  \\
\hline
\end{tabular}

\vspace{0.8cm}
\noindent\textbf{Section B: Fill in the Blanks (1--5)}\\[0.2cm]
\noindent
1.\ \underline{\hspace{17cm}}\\[0.4cm]
2.\ \underline{\hspace{17cm}}\\[0.4cm]
3.\ \underline{\hspace{17cm}}\\[0.4cm]
4.\ \underline{\hspace{17cm}}\\[0.4cm]
5.\ \underline{\hspace{17cm}}\\[0.4cm]

\vspace{0.6cm}
\noindent\textbf{Section C: Short Answers (1--5)}\\[0.2cm]
\noindent
1.\ \underline{\hspace{17cm}}\\[0.35cm]
\underline{\hspace{17cm}}\\[0.5cm]
\underline{\hspace{17cm}}\\[0.5cm]
\underline{\hspace{17cm}}\\[0.5cm]
2.\ \underline{\hspace{17cm}}\\[0.35cm]
\underline{\hspace{17cm}}\\[0.5cm]
\underline{\hspace{17cm}}\\[0.5cm]
\underline{\hspace{17cm}}\\[0.5cm]
3.\ \underline{\hspace{17cm}}\\[0.35cm]
\underline{\hspace{17cm}}\\[0.5cm]
\underline{\hspace{17cm}}\\[0.5cm]
\underline{\hspace{17cm}}\\[0.5cm]
4.\ \underline{\hspace{17cm}}\\[0.35cm]
\underline{\hspace{17cm}}\\[0.5cm]
\underline{\hspace{17cm}}\\[0.5cm]
\underline{\hspace{17cm}}\\[0.5cm]
5.\ \underline{\hspace{17cm}}\\[0.35cm]
\underline{\hspace{17cm}}\\[0.5cm]
\underline{\hspace{17cm}}\\[0.5cm]
\underline{\hspace{17cm}}\\[0.5cm]
\vspace{0.6cm}
\noindent\textbf{Section D: Design Questions (1--2)}\\[0.2cm]
\noindent
\textbf{D1:}\\[0.2cm]
    \underline{\hspace{17cm}}\\[0.35cm]
\underline{\hspace{17cm}}\\[0.35cm]
\underline{\hspace{17cm}}\\[0.35cm]
\underline{\hspace{17cm}}\\[0.35cm]
\underline{\hspace{17cm}}\\[0.35cm]
\underline{\hspace{17cm}}\\[0.35cm]
\underline{\hspace{17cm}}\\[0.35cm]
\underline{\hspace{17cm}}\\[0.35cm]
\underline{\hspace{17cm}}\\[0.35cm]

\newpage
\textbf{D2:}\\[0.2cm]
\underline{\hspace{17cm}}\\[0.35cm]
\underline{\hspace{17cm}}\\[0.35cm]
\underline{\hspace{17cm}}\\[0.35cm]
\underline{\hspace{17cm}}\\[0.35cm]
\underline{\hspace{17cm}}\\[0.35cm]
\underline{\hspace{17cm}}\\[0.35cm]
\underline{\hspace{17cm}}\\[0.35cm]
\underline{\hspace{17cm}}\\[0.35cm]
\underline{\hspace{17cm}}\\[0.35cm]

\end{document}
